\chapter{尽职调查：财务、法律与商业}

\section{DD 的整体节奏}

尽职调查（Due Diligence）应在 Term Sheet 签署后、重大交割款项划出前完成，但\textbf{高风险事项}（诉讼、环保、核心知识产权归属）宜在签署前做\textbf{有限度预尽调}以避免排他期浪费。典型工作流：数据室清单 → 管理层问答（Q\&A）→ 现场走访 → 第三方报告（环境、税务、IT）→ 问题汇总与价格/条款调整建议。投资团队须指定\textbf{工作流负责人}，防止律师与会计师各跑各的、结论无法合并。

\section{财务尽调}

关注\textbf{盈利质量}：收入确认政策是否激进？应收账款与现金流是否匹配？一次性收益与非经常性项目是否被误当作可持续利润？\textbf{正常化 EBITDA} 的调整须有证据支撑并在投资备忘录中披露假设。营运资金（NWC）目标与交割后「漏桶」条款常是谈判焦点。对早期企业，财务 DD 可能让位于\textbf{单位经济与跑道}分析。

\section{法律尽调}

重点扫描：重大合同（客户/供应商集中度、变更控制权条款）、诉讼与监管调查、土地与环保许可、知识产权链条（职务发明、开源许可污染）、劳动与社保合规、数据隐私。红色事项应进入\textbf{交割条件（CP）}或价格机制（earn-out、托管账户），而非仅停留在尽调报告附录。

\section{商业与运营尽调}

访谈中层与客户样本（在合规前提下）、考察供应链与产能利用率、评估数字化与内控成熟度。对 PE 而言，\textbf{100 日计划}的雏形往往在此阶段形成：投后哪些杠杆（采购、定价、组织）最可能兑现。

\section{内部与外部顾问的协作}

内部：投资负责人对结论负总责，需参与关键会议而非仅读摘要。外部：律师与审计师的范围说明书（SOW）应明确交付物与时间，避免范围蔓延。所有重大发现应进入\textbf{问题清单}并标注状态（已解决/留待交割后/放弃交易）。
